% =====================================================================
%  Capitolo 1 — Introduzione
% =====================================================================
\chapter{Introduzione}
\label{cap:introduzione}

\section{Contesto}
\label{sec:contesto}

Le immagini telerilevate ad altissima risoluzione (\emph{Very High Resolution}, VHR), riprese
da aereo o da satellite con una distanza di campionamento al suolo (\emph{Ground Sampling
Distance}, GSD) di pochi centimetri, sono oggi una fonte primaria di informazione per il
monitoraggio delle aree urbane, la pianificazione dell'uso del suolo, la tutela ambientale e la
gestione del traffico~\cite{zhu2017dlrs,cheng2016survey}. Per trasformare queste immagini in
informazione utilizzabile occorre sapere \emph{che cosa} c'è in ogni punto della scena: è il
compito della \emph{segmentazione semantica}, che assegna a ciascun pixel una classe scelta da un
insieme fissato a priori.

Rispetto alla classificazione di un'immagine intera o al rilevamento di oggetti con riquadri, la
segmentazione semantica è una predizione \emph{densa}: per un'ortofoto di $6000\times6000$ pixel
il modello produce $3{,}6\cdot10^{7}$ decisioni, che devono rispettare i contorni degli oggetti e
rimanere coerenti su scala urbana. Dal 2015 il problema è affrontato con reti neurali
convoluzionali completamente convoluzionali~\cite{long2015fcn} e con architetture
encoder-decoder come U-Net~\cite{ronneberger2015unet}. Negli anni successivi queste architetture
sono state arricchite con convoluzioni dilatate~\cite{chen2018deeplab}, meccanismi di
attenzione~\cite{wang2018nonlocal,fu2019danet} e, più di recente, con encoder di tipo
transformer~\cite{dosovitskiy2021vit,liu2021swin,xie2021segformer}.

Il banco di prova di riferimento per la segmentazione di scene urbane aeree è il \emph{2D
Semantic Labeling Contest} dell'ISPRS~\cite{rottensteiner2012isprs,isprs2dlabeling}, con i
dataset \textbf{Potsdam} (ortofoto RGB a \SI{5}{\centi\meter}) e \textbf{Vaihingen} (falso colore
IR-R-G a \SI{9}{\centi\meter}), annotati in sei classi: superfici impermeabili
\ClasseBox{ClsImpervious}, edifici \ClasseBox{ClsBuilding}, vegetazione bassa
\ClasseBox{ClsLowVeg}, alberi \ClasseBox{ClsTree}, auto \ClasseBox{ClsCar} e clutter
\ClasseBox{ClsClutter}.

\section{Le difficoltà del dominio}
\label{sec:difficolta}

Il telerilevamento urbano VHR combina quattro difficoltà, che motivano le scelte di questa tesi
(\cref{fig:filo-logico}).

\begin{description}[style=nextline, leftmargin=1.2em]
  \item[Variabilità intra-classe alta e inter-classe bassa.] Gli edifici hanno materiali,
  forme e colori molto diversi tra loro, mentre tetti, marciapiedi e strade in cemento possono
  essere quasi indistinguibili nel solo spettro visibile. Lo stesso vale per alberi e
  vegetazione bassa, che si distinguono soprattutto per l'altezza.
  \item[Condizioni di ripresa.] Ombre degli edifici alti, foschia e differenze di
  illuminazione alterano l'aspetto delle superfici e introducono confusioni sistematiche.
  \item[Scarsità di etichette.] A differenza dei grandi dataset di immagini naturali, le
  annotazioni pixel per pixel di immagini aeree sono poche, costose e richiedono esperti: due
  delle collezioni ISPRS più usate contano 38 (Potsdam) e 33 (Vaihingen) aree annotate.
  \item[Contesto ampio.] Classificare un pixel richiede di guardare l'oggetto e il suo intorno:
  un'auto si riconosce dalla strada su cui si trova, un cortile dagli edifici che lo
  circondano. Aggregare contesto su immagini grandi è però costoso. L'attenzione
  \emph{dot-product} ha un costo quadratico nel numero di posizioni, già proibitivo su un
  singolo ritaglio di $512\times512$ pixel (\cref{sec:attenzione-dot}).
\end{description}

\begin{figure}[t]
  \centering
  \resizebox{\linewidth}{!}{% Schema: dai problemi del dominio ai meccanismi studiati e alle domande di ricerca.
\begin{tikzpicture}[
    node distance=4mm and 9mm,
    prob/.style={dato, text width=3.5cm, minimum height=10mm},
    mecc/.style={attn, text width=3.4cm, minimum height=10mm},
    domanda/.style={blocco, fill=white, draw=SapRosso, text width=2.6cm, minimum height=10mm},
  ]
  % Colonna 1: problemi
  \node[prob] (p1) {Immagini molto grandi\\{\scriptsize $6000\times6000$ px, contesto ampio}};
  \node[prob, below=of p1] (p2) {Variabilit\`a intra-classe alta,\\inter-classe bassa};
  \node[prob, below=of p2] (p3) {Etichette scarse e costose};
  \node[prob, below=of p3] (p4) {Ambiguit\`a spettrale\\{\scriptsize tetto vs strada, albero vs prato}};

  % Colonna 2: meccanismi
  \node[mecc, right=of p1] (m1) {Attenzione lineare\\{\scriptsize LAM, costo $\mathcal{O}(N)$}};
  \node[mecc, right=of p2] (m2) {Bottleneck discreto\\{\scriptsize codebook VQ}};
  \node[mecc, right=of p3] (m3) {Auto-distillazione\\{\scriptsize DINO su dati non etichettati}};
  \node[mecc, right=of p4] (m4) {Altezza sul terreno\\{\scriptsize nDSM come 4\textordmasculine{} canale}};

  % Colonna 3: domande di ricerca
  \node[domanda, right=of m1] (r1) {\RQ{1}, \RQ{5}\\{\scriptsize contributo e costo}};
  \node[domanda, right=of m2] (r2) {\RQ{2}\\{\scriptsize codebook sano?}};
  \node[domanda, right=of m3] (r3) {\RQ{3}, \RQ{4}\\{\scriptsize poche etichette, augmentation}};
  \node[domanda, right=of m4] (r4) {\RQ{6}\\{\scriptsize fusione precoce}};

  \foreach \i in {1,...,4}{
    \draw[flusso] (p\i) -- (m\i);
    \draw[flusso] (m\i) -- (r\i);
  }
  % legami secondari
  \draw[flusso, densely dotted] (p2.east) -- (m4.north west);
  \draw[flusso, densely dotted] (p3.east) -- (m2.south west);

  % Intestazioni
  \node[nota, above=2mm of p1, font=\footnotesize\bfseries, text=SapRossoScuro] {Problema};
  \node[nota, above=2mm of m1, font=\footnotesize\bfseries, text=SapRossoScuro] {Meccanismo};
  \node[nota, above=2mm of r1, font=\footnotesize\bfseries, text=SapRossoScuro] {Domanda};

  % Riquadro: estensioni rispetto al 2022
  \begin{scope}[on background layer]
    \node[gruppo, fit=(m4)(r4), label={[nota]below:estensione rispetto alla tesi del 2022}] {};
  \end{scope}
\end{tikzpicture}
}
  \caption[Dai problemi del dominio alle domande di ricerca]{Dai problemi del dominio ai
    meccanismi studiati e alle domande di ricerca. Le linee punteggiate indicano legami
    secondari: l'altezza aiuta a separare classi spettralmente simili e la quantizzazione nasce
    dall'idea di \emph{visual words} usata dal pre-addestramento OBoW.}
  \label{fig:filo-logico}
\end{figure}

\section{Dalla tesi del 2022 a questo lavoro}
\label{sec:da-2022}

Questa tesi riprende il lavoro svolto nel 2022~\cite{campana2022tesi}, a sua volta continuazione
di MARE~\cite{marsocci2021mare}. MARE ha mostrato che l'architettura
\maresunet{}~\cite{li2022maresunet}, cioè un encoder ResNet~\cite{he2016resnet} e un decoder
U-Net con un \emph{Linear Attention Mechanism}~(LAM)~\cite{li2020lam} nelle connessioni di
salto, beneficia sul dataset Vaihingen di un pre-addestramento auto-supervisionato basato su
\emph{bag of visual words}~(OBoW)~\cite{gidaris2021obow}. La tesi del 2022 ha esteso
quell'idea in tre direzioni: un dataset più grande (Potsdam), un metodo di auto-distillazione
più recente (DINO~\cite{caron2021dino}) e un collo di bottiglia esplicitamente discreto, con un
codebook di quantizzazione vettoriale~\cite{vandenoord2017vqvae} ispirato al vocabolario di
\emph{visual words} di OBoW.

Il risultato sperimentale di allora è stato negativo per entrambe le estensioni: la
configurazione migliore restava la \maresunet{} con encoder ResNet-18 pre-addestrato su
ImageNet (mIoU di test $\approx 0{,}604$). Una revisione sistematica del protocollo e del codice
originale (\cref{sec:lettura-critica}) mostra tuttavia che quel risultato \emph{non è
conclusivo}. La classe clutter non riceveva gradiente, le metriche erano mediate per batch, la
validation probabilmente non era indipendente dal training e il pre-addestramento si
discostava da DINO in punti sostanziali. Le domande poste nel 2022 restano quindi aperte e
meritano una risposta basata su un protocollo corretto.

\section{Obiettivi}
\label{sec:obiettivi}

L'obiettivo della tesi è \textbf{rivalutare in modo riproducibile e confrontabile con la
letteratura} il contributo di tre meccanismi (attenzione lineare, quantizzazione vettoriale e
auto-distillazione) alla segmentazione semantica di immagini aeree urbane, ed estendere
l'analisi all'uso dell'informazione altimetrica. In dettaglio:

\begin{enumerate}
  \item formalizzare i meccanismi studiati e derivarne le proprietà rilevanti per il dominio:
  costo computazionale, capacità informativa del collo di bottiglia, modi di collasso
  dell'auto-distillazione (\cref{cap:teoria});
  \item definire un protocollo sperimentale che elimini i fattori confondenti individuati nel
  lavoro del 2022 (\cref{cap:metodologia});
  \item rispondere con più seed e intervalli di variabilità alle domande di ricerca della
  \cref{sec:domande} (\cref{cap:esperimenti});
  \item quantificare il compromesso tra accuratezza e costo su hardware consumer (una GPU da
  \SI{8}{\giga\byte}), in linea con la motivazione energetica già presente nel 2022.
\end{enumerate}

\section{Domande di ricerca}
\label{sec:domande}

\begin{description}[style=nextline, leftmargin=1.2em]
  \item[\RQ{1} -- Attenzione lineare.] Su Potsdam, con un protocollo standard, quanto rende
  \maresunet{} rispetto a una rete con la stessa struttura ma connessioni di salto senza
  attenzione, a parità di encoder?
  \item[\RQ{2} -- Quantizzazione vettoriale.] Un collo di bottiglia quantizzato, con un codebook
  \emph{sano} (perplexity alta, pochi codici inutilizzati), migliora la segmentazione? Conta dove
  la quantizzazione viene applicata: solo al collo di bottiglia o anche alle connessioni di
  salto?
  \item[\RQ{3} -- Auto-distillazione e poche etichette.] Un pre-addestramento DINO fedele
  all'originale, su immagini di dominio non etichettate, eguaglia o supera ImageNet? Come cambia
  il confronto quando si usa solo l'1\%, il 10\% o il 25\% delle etichette?
  \item[\RQ{4} -- Augmentation per immagini aeree.] Le trasformazioni cromatiche
  \enquote{forti} di DINO sono adatte alle immagini aeree, o conviene una variante attenuata e
  invariante alle rotazioni?
  \item[\RQ{5} -- Efficienza.] Qual è il compromesso tra accuratezza e costo (parametri,
  memoria, throughput) tra encoder ResNet-18, ResNet-34 e ResNet-50 e tra encoder
  convoluzionali e transformer?
  \item[\RQ{6} -- Altezza sul terreno.] L'nDSM come quarto canale di ingresso migliora la
  segmentazione, in particolare per le coppie edifici/superfici impermeabili e alberi/vegetazione
  bassa? Il risultato cambia se l'nDSM è quello ufficiale o quello stimato dal DSM?
\end{description}

Le domande da \RQ{1} a \RQ{5} riformulano le ipotesi del 2022. La \RQ{6} e il secondo benchmark
(Vaihingen) sono estensioni introdotte in questo lavoro.

\section{Contributi}
\label{sec:contributi}

\begin{itemize}
  \item Una \textbf{lettura critica} della tesi del 2022 che separa le conclusioni giustificate
  dai dati da quelle dovute a difetti di protocollo o di implementazione (\cref{cap:premesse}).
  \item Una \textbf{trattazione teorica unificata} dei meccanismi studiati. Comprende la
  derivazione del LAM come attenzione a nucleo con costo $\bigO{N}$ e una sua proprietà di
  \enquote{basso contrasto}, l'interpretazione dell'aggiornamento EMA del codebook come K-means
  online, l'analisi dell'aggiramento del collo di bottiglia da parte delle connessioni di salto,
  i modi di collasso dell'auto-distillazione e l'effetto di una classe esclusa dalla perdita
  (\cref{cap:teoria}).
  \item Un \textbf{protocollo sperimentale} con split standard per tile intere, valutazione su
  tile complete con etichette piene ed \emph{eroded}, matrice di confusione globale, più seed e
  un regime a poche etichette con budget di iterazioni costante (\cref{cap:metodologia}).
  \item Un'\textbf{implementazione nuova e verificata} (pacchetto \rsseg{}, Keras~3 su backend
  PyTorch, tracciamento con MLflow). Comprende campagne di esperimenti riprendibili, 109 test
  automatici e la verifica di equivalenza numerica con i modelli di riferimento
  (\cref{sec:implementazione}).
  \item I \textbf{risultati sperimentali} che rispondono alle domande di ricerca
  (\cref{cap:esperimenti}).
  \Segnaposto{Riassumere i risultati principali in una riga per domanda di ricerca, quando
  disponibili.}
\end{itemize}

\section{Struttura della tesi}
\label{sec:struttura}

Il \cref{cap:premesse} inquadra il lavoro nello stato dell'arte e analizza criticamente la tesi
del 2022. Il \cref{cap:teoria} sviluppa il quadro teorico: formulazione del problema, encoder-decoder
residuali, attenzione lineare, quantizzazione vettoriale, auto-distillazione, fusione
dell'nDSM, metriche e inferenza su tile intere. Il \cref{cap:metodologia} descrive dati,
protocolli, configurazioni di addestramento, implementazione e analisi statistica. Il
\cref{cap:esperimenti} è dedicato agli esperimenti, organizzati per domanda di ricerca. Il
\cref{cap:prospettive} discute lo stato del lavoro, le attività rimanenti e le prospettive di
ricerca.
